% Options for packages loaded elsewhere
\PassOptionsToPackage{unicode}{hyperref}
\PassOptionsToPackage{hyphens}{url}
%
\documentclass[
]{article}
\usepackage{amsmath,amssymb}
\usepackage{lmodern}
\usepackage{iftex}
\ifPDFTeX
  \usepackage[T1]{fontenc}
  \usepackage[utf8]{inputenc}
  \usepackage{textcomp} % provide euro and other symbols
\else % if luatex or xetex
  \usepackage{unicode-math}
  \defaultfontfeatures{Scale=MatchLowercase}
  \defaultfontfeatures[\rmfamily]{Ligatures=TeX,Scale=1}
\fi
% Use upquote if available, for straight quotes in verbatim environments
\IfFileExists{upquote.sty}{\usepackage{upquote}}{}
\IfFileExists{microtype.sty}{% use microtype if available
  \usepackage[]{microtype}
  \UseMicrotypeSet[protrusion]{basicmath} % disable protrusion for tt fonts
}{}
\makeatletter
\@ifundefined{KOMAClassName}{% if non-KOMA class
  \IfFileExists{parskip.sty}{%
    \usepackage{parskip}
  }{% else
    \setlength{\parindent}{0pt}
    \setlength{\parskip}{6pt plus 2pt minus 1pt}}
}{% if KOMA class
  \KOMAoptions{parskip=half}}
\makeatother
\usepackage{xcolor}
\setlength{\emergencystretch}{3em} % prevent overfull lines
\providecommand{\tightlist}{%
  \setlength{\itemsep}{0pt}\setlength{\parskip}{0pt}}
\setcounter{secnumdepth}{-\maxdimen} % remove section numbering
\ifLuaTeX
  \usepackage{selnolig}  % disable illegal ligatures
\fi
\IfFileExists{bookmark.sty}{\usepackage{bookmark}}{\usepackage{hyperref}}
\IfFileExists{xurl.sty}{\usepackage{xurl}}{} % add URL line breaks if available
\urlstyle{same} % disable monospaced font for URLs
\hypersetup{
  hidelinks,
  pdfcreator={LaTeX via pandoc}}

\author{Dietmar Gerald Schrausser}
\date{09.05.2026}


\begin{document}

\hypertarget{foliumblat-iiiir}{%
\section{Foliu{[}m{]}/Blat IIIIr}\label{foliumblat-iiiir}}

A \emph{contemporary} English translation is presented from comparing
the existing Latin, Early New High German (ENHG) and English
(translation) texts for better comprehensibility. Especially since the
ENHG text is difficult to understand compared to modern High German, as
a \emph{profound} knowledge of German (as a mother tongue) as well as
already aknowledge of various \emph{dialects} is required.

\hypertarget{on-the-work-of-the-fourth-day}{%
\subsection{On the Work of the Fourth
Day}\label{on-the-work-of-the-fourth-day}}

\begin{quote}
\textbf{Q}Uarto die dixit deus. Fia{[}n{]}t luminaria in
firmame{[}n{]}to celi: {[}et{]} diuidant diem et nocte{[}m{]}.\\
Et sint in signa {[}et{]} t{[}em{]}p{[}er{]}a: {[}et{]} dies: {[}et{]}
a{[}n{]}nos: vt lucea{[}n{]}t in firmame{[}n{]}to celi: {[}et{]}
illumine{[}n{]}t terra{[}m{]}.\\
Et factu{[}m{]} e{[}st{]} ita.\\
Fecitq{[}ue{]} deus duos luminaria magna: luminare maius: vt
p{[}re{]}esset diei: {[}et{]} luminare min{[}us{]}: vt p{[}er{]}esset
nocti {[}et{]} stellas: vt duiderent luce{[}m{]} {[}et{]} tenebras.
\end{quote}

\textbf{Moses} first mentions the celestial bodies\footnote{`celestium'
  und `himlischen ding'.} that God placed in the firmament so that they
might shine in the heaven\footnote{`celo' und `himel'.} and illuminate
the \emph{earth}, namely the sun, moon, and stars\footnote{`stellas' und
  `stern'.} with which the upper\footnote{`superior' und `oberteil'.}
part of the \emph{world} itself is adorned, just as the \emph{earth} is
adorned with the metals, plants and living beings\footnote{`animantibus'.}
that grow in it. After speaking about the nature of the firmament, it
remained for him to speak about the workings of the stars\footnote{`siderum'
  und `gestirns'.} and their function, to explain for what purpose they
were created and for what task God assigned them.\footnote{`Celestium
  enim corporum' und `himlischen leiplichen ding'.}

The two modes of action of celestial bodies\footnote{from the
  `Hexameron' (Ambrosius,
  \href{https://books.google.com/books?id=L1DbnAAACAAJ}{1897}, Book 4,
  Chapter 1).}, \emph{movement} and \emph{illumination}, are universally
recognizable, resulting in a \emph{twofold} movement (I): One is (i) the
entire world\footnote{`mundi totius' und `ga{[}n{]}tze{[}n{]} werlt'.},
in which heaven and ether rotate in a perfect circle through the entire
space of the universe\footnote{`vniuersi' und `werlt'.} in 24 hours. The
other is the actual, multifaceted, and varied movement of (ii) the
celestial bodies\footnote{`siderum' und `gestirns'.}, with the main
movement of the sun, which passes through all the signs of the zodiac
within twelve months. This creates the day, from which the daily cycles
form the year. The remaining movements of the stars occur at different
intervals. Rightly, \emph{Moses} aptly and succinctly reminded us that
the stars were placed in the firmament for days, years, and
seasons.\footnote{c.f. `Heptaplus' (Mirandola \& Mirandola,
  \href{https://doi.org/10.3931/e-rara-61217}{1601}).}

He also explicitly pointed out the other effect of the stars, which is
\emph{illumination}\footnote{`illuminatio' und `erlewchtu{[}n{]}g'.}
(II), when he said that they were arranged to shine in the
heaven\footnote{`vt lucerent in celo' und `zu scheine{[}n{]} am{[}m{]}
  hymel'.} and to light up the earth. Thus, the bodies of the moon, the
sun and the stars\footnote{`stellarum' und `stern{[}n{]}'.} are
distributed\footnote{`distributa' und `außgetailt'.} for this
purpose\footnote{`ministeria' und `dienstperkeiten'.}. Although the sun
only \emph{apparently} rises alone during the day (\textbf{Cicero}
probably wants it to appear that only \emph{one} sun is
recognized)\footnote{Refers to \emph{De Natura Deorum} (Cicero \&
  Cicero, \href{https://doi.org/10.3931/e-rara-89116}{1496}, 2.68),
  where it is suggested that the sun is called \emph{sol} (solus)
  because it is the only one of its kind, or because it \emph{appears}
  alone after it has obscured all the other stars.}, since it
\emph{appears} alone through the veiled stars\footnote{`sideribus'.},
\emph{it} (the sun) is nevertheless the true and perfect light that
illuminates everything with the greatest \emph{warmth} and the brightest
\emph{brilliance}\footnote{`fulgore clarissimo' und `allerklarste{[}m{]}
  schein'.}. For although countless stars seem to twinkle and shine,
those whose light is not full and clear neither provide warmth nor
dispel the darkness through their multitude. Thus, \emph{two} main
elements can be found, (i) \emph{warmth} and (ii)
\emph{`glory'}\footnote{`humor' und `feüchtigkeit' (to dispel the
  darkness).}, which possess a diverse and at the same time cooperating
power, wonderfully created by God for the preservation and creation of
all things.\footnote{from the `Speculum Naturale' (De Beauvais,
  \href{https://nbn-resolving.org/urn:nbn:de:hbz:061:1-20612}{1964},
  Buch 1, Kapitel 75).}

However, this would require addressing \emph{very profound} questions,
each of which would require its own book: (i) In what way are these
stars located in the firmament, whether as its \emph{nobler} parts or as
animals\footnote{`a{[}n{]}{[}m{]}alia'(sic) und `geschöpff'.} in their
spheres, like fish in water or animals on earth.\footnote{This reflects
  a classic philosophical debate of the Middle Ages: are celestial
  bodies merely physical components of the heaven, or do they possess
  their own \emph{soul}?} This point also required a discussion with
(ii) astrologers\footnote{`Genethliacis'.} regarding divination by the
stars and the prediction of future events, which science confirmed by
stating that \emph{Moses} said the stars were placed by God as signs.
Furthermore, (iii) the nature of the stars, their motion, their orbits,
the \emph{spots}\footnote{`maculis'.} on the moon, and the entire
science of the stars (astronomy \emph{and} astrology)\footnote{`siderali
  scie{[}n{]}tia'.} must be addressed.\footnote{from the `Heptaplus'
  (Mirandola \& Mirandola,
  \href{https://doi.org/10.3931/e-rara-61217}{1601}).}

Truly,

\begin{quote}
qua{[}m{]}q{[}uoniam{]} sint pulcra {[}et{]} digna cognitu.\footnote{`Although
  they are beautiful and deserve recognition' is frequently used to
  describe the works of Horace
  (\href{https://google.cat/books?id=A28VAAAAYAAJ}{1868}), summarizing
  his famous principle from the \emph{Ars Poetica}.}
\end{quote}

Maybe we will hear \emph{this} \textbf{Horace}, but here was not the
right place.\footnote{`sed nunc non erat hic locus', Horace
  (\href{https://google.cat/books?id=A28VAAAAYAAJ}{1868}, \emph{Ars
  Poetica}, Lines 14--19).}

\hypertarget{references}{%
\subsection{References}\label{references}}

Ambrosius, A. (1897). \emph{Opera: Exameron. De Paradiso. De Cain Et
Abel. De Noe. De Abraham. De Isaac. De Bono Mortis}. Edited by Faller,
O., \& Schenkl, K. Corpus Scriptorum Ecclesiasticorum Latinorum.
Hoelder-Pichler-Tempsky.
\url{https://books.google.com/books?id=L1DbnAAACAAJ}

Cicero, M. T., \& Cicero, Q. (1496). \emph{M.T. Ciceronis de Natura
Deorum Libri Tres} \ldots{} Impræssum Venetiis: per Symonem Papiensem
dictum Biuilaqua. \url{https://doi.org/10.3931/e-rara-89116}

De Beauvais, V. (1964). \emph{Speculum maius}. Graz: Akad. Druck- u.
Verl.-Anst. \url{https://nbn-resolving.org/urn:nbn:de:hbz:061:1-20612}

Horace. (1868). \emph{Q. Horatii Flacci Opera Omnia}. Oxoni et Londini:
J. Parker et Soc. \url{https://google.cat/books?id=A28VAAAAYAAJ}

Lactantius, L.C.F. (1497). \emph{Divinae Institutiones}. Venetiis: Simon
Bevilaqua. \url{https://books.google.com/books?id=qGsnghsi3uUC}

Mirandola, G., \& Mirandola, G. F. (1601). \emph{Ioannis Pici,
Mirandulae \ldots{} opera quae extant omnia \ldots{}} Basileae: per
Sebastianum Henricpetri. \url{https://doi.org/10.3931/e-rara-61217}

\end{document}
